\textbf{Lösung zu Übung 4, Aufgabe 8}

\begin{loesung}
% KORRIGIERT: Kontext "in Z_17" für die Gleichungen ergänzt
Alle Rechnungen in $\Zm{17}$:

\begin{tabular}{lcrcr}
$2,\ 4,\ 8,\ -1,\ -2,\ -4,\ -8,\ 1$ & $\Rightarrow$ & $\operatorname{ord}(2) = 8$ & $\Rightarrow$ & $2^{-1} = -8$\\
$3,\ -8,\ -7,\ -4,\ 5,\ -2,\ -6,\ -1,\ \ldots$ & $\Rightarrow$ & $\operatorname{ord}(3) = 16$ & $\Rightarrow$ & $3^{-1} = 6$\\
$4,\ -1,\ -4,\ 1$ & $\Rightarrow$ & $\operatorname{ord}(4) = 4$ & $\Rightarrow$ & $4^{-1} = -4$\\
$5,\ 8,\ 6,\ -4,\ -3,\ 2,\ -7,\ -1,\ \ldots$ & $\Rightarrow$ & $\operatorname{ord}(5) = 16$ & $\Rightarrow$ & $5^{-1} = 7$
\end{tabular}

\medskip
Primitiv sind 3, 5
\begin{align*}
2 \cdot 9 &= 18 \equiv 1 \pmod{17}\\
3 \cdot 6 &= 18 \equiv 1 \pmod{17}\\
4 \cdot 13 &= 52 \equiv 1 \pmod{17}\\
5 \cdot 7 &= 35 \equiv 1 \pmod{17}
\end{align*}
\end{loesung}
